\documentclass[11pt,a4paper]{article}
\usepackage{fontspec,amsmath,amssymb,amsthm,microtype}
\setmainfont{lmroman10-regular.otf}[BoldFont=lmroman10-bold.otf,ItalicFont=lmroman10-italic.otf,BoldItalicFont=lmroman10-bolditalic.otf]
\usepackage[margin=25mm]{geometry}
\usepackage[unicode,hidelinks]{hyperref}
\hypersetup{pdfauthor={LangevinLean Project}}
\setlength{\emergencystretch}{3em}
\setlength{\parskip}{0.4em}
\setlength{\parindent}{0pt}
\allowdisplaybreaks[2]
\raggedbottom
\title{\texorpdfstring{Derive the actual stationary score-\allowbreak{}increment pairing under the printed C3 condition}{Derive the actual stationary score-increment pairing under the printed C3 condition}}
\author{LangevinLean Project}
\date{}
\begin{document}\maketitle
Draft mathematical exposition; these notes may contain mistakes.


This module supplies the stochastic pairing estimate needed in the Laplacian one-step proof. It starts from the actual Brownian Langevin integral equation, Gaussian-innovation independence, and Gibbs stationarity. Compact C2 score approximations satisfy the proved weighted Dynkin identity. Their genuine generator errors have a stationary L2 norm that is uniform in time. Taking the compact approximation limit and then the initial-weight clipping limit yields the sharp coefficient \(M_\infty\) sqrt(p \(M_{\rm av}\))+\(M_\Delta\) sqrt(p). Only the source C3 potential and its bound on gradient(Laplacian f) are used.



Stationarity transfers every genuine Gibbs squared approximation error to the actual observation:


\[\|F_n(Y_t)-\nabla f(Y_t)\|_{L^2(P)}^2
=\int\|F_n(x)-\nabla f(x)\|^2\,\pi(dx)\longrightarrow0.\]


Applying the compact Dynkin identities, the generator approximation, and radial clipping proves the sharp cross-term bound used in Theorem 2:


\[\left\langle\Psi,\nabla f(Y_T)-\nabla f(Y_0)\right\rangle_{L^2(P)}
\le T\|\Psi\|_{L^2(P)}
\left(M_\infty\sqrt{M_{\mathrm{av}}p}+M_\Delta\sqrt p\right).\]
\section{\texorpdfstring{Transfer Gibbs squared-\allowbreak{}error convergence through an actual stationary law}{Transfer Gibbs squared-error convergence through an actual stationary law}}


If a random L2 state has law pi, the squared norm of the difference of its two canonical observations is exactly the integral of the squared pointwise error under pi. Taking square roots identifies the L2 distance. Continuity of the square root at zero therefore turns convergence of the genuine squared-error integrals into convergence of the actual L2 observations. This theorem uses the displayed map-law equality and the square-integrability witnesses, rather than assuming convergence after composition.


\section{\texorpdfstring{Construct and pass the actual compact Dynkin pairings to the score increment}{Construct and pass the actual compact Dynkin pairings to the score increment}}


Strong convexity normalizes the Gibbs measure and gives its finite partition. The integral equation at time zero identifies the initial state with Y(0), so stationarity supplies its Gibbs law as well as that of the final state. The proved radial score-cutoff construction provides compact C2 scalar tests in every coordinate, with explicit L2 membership of their vector tests and actual Langevin generators. Their squared Gibbs errors tend to zero against the actual gradient and gradient generator. The observation convergence theorem yields both endpoint L2 limits.



The gradient generator is the actual vector gradient(Laplacian f) minus Hessian f applied to gradient f. Its stationary L2 norm is at most R=\(M_\infty\) sqrt(p \(M_{\rm av}\))+\(M_\Delta\) sqrt(p), by the exact average-smoothness Gibbs score bound and the printed Laplacian-gradient condition. At every time the observation of an approximating generator differs from this target by the same Gibbs L2 error err(n), independent of time. The triangle inequality therefore bounds every approximating generator norm by R+err(n), and err(n) tends to zero. The cutoff construction needs only C2 regularity of the score, which follows from C3 regularity of f.



Radially clip the initial vector Psi. The clipped vector has bounded coordinates and remains independent of future Gaussian innovations jointly with the initial state, because measurable maps preserve independence. Apply the actual compact-test weighted Dynkin theorem to every scalar coordinate weight and its compact C2 score approximation. Compact generators are bounded and uniformly continuous, so these weighted generator paths are interval integrable. Sum the coordinate identities, exchange only the finite sum with the time integral, and identify the sum with the exact Euclidean L2 pairing.



For a fixed clipped weight, apply the compact Dynkin approximation limit theorem to the endpoint L2 limits and the uniform generator bound R+err(n). This gives the pairing of the clipped weight with the actual score increment at most T times its L2 norm times R. Finally apply the clipping limit theorem: the pairings converge, and the clipped norms never exceed the original L2 norm. The resulting stationary gradient pairing bound applies to every square-integrable independent initial vector. Its stochastic hypotheses concern the actual Brownian integral equation and independence; no desired Wasserstein recurrence or score-increment bound is supplied as a premise.



Earlier results used here: \cite[Section~3.13.2]{LangevinLean2026ULASC}; \cite[Section~3.13.3]{LangevinLean2026ULASC}; \cite[Section~3.25.1]{LangevinLean2026ULASC}; \cite[Section~3.25.3]{LangevinLean2026ULASC}; \cite[Section~3.26.10]{LangevinLean2026ULASC}; \cite[Section~3.26.13]{LangevinLean2026ULASC}; \cite[Section~3.26.3]{LangevinLean2026ULASC}; \cite[Section~3.29.7]{LangevinLean2026ULASC}; \cite[Section~3.46.2]{LangevinLean2026ULASC}.

\begin{thebibliography}{ULA-LSI}
\bibitem[ULA-SC]{LangevinLean2026ULASC}
LangevinLean Project. \emph{Unadjusted Langevin under Strong Convexity}. LL-01, 2026. Draft report, content revision 6f6f50decdb6. \url{https://odakyildiz.github.io/LangevinLean/docs/pdf/UnadjustedLangevin/FullProof.pdf}.
\end{thebibliography}

\end{document}
